
\chapter{Conclusion and Perspectives}
\label{chap:conclusion}

This project set out to build a neural architecture capable of learning and advancing the time evolution of physical systems governed by partial differential equations. Starting from the Neural Field Turing Machine (NFTM) framework~\cite{malhotra2025neuralfieldturingmachine}, we investigated a progression of controllers (CNN, Transformer, and TCAN) on the 1D viscous Burgers equation, extended the best-performing model to 2D, and benchmarked it against the Fourier Neural Operator~\cite{li2021fourier}. This chapter summarises our main findings, the limitations we encountered, and the perspectives for future work.

%==============================================================================
\section{Summary of Contributions}
\label{sec:conclusion_summary}
%==============================================================================

\paragraph{Controller comparison.}
We evaluated three controller architectures within the NFTM framework on long-horizon autoregressive rollouts of 1D Burgers dynamics. The CNN controller excels at exploiting local spatial patterns but struggles to capture long-range temporal dependencies, leading to error accumulation over extended rollouts. The Transformer controller handles arbitrary temporal receptive fields, but its cost grows with sequence length and it requires more parameters to match the accuracy of the TCAN. The \textbf{Causal Temporal Convolutional Attention Network (TCAN)} achieves the best results by combining spatial convolutional feature extraction with a causal attention mechanism that also aggregates the full history window. This design, the best of the CNN and Transformer worlds, is the principal model we adopt for all quantitative evaluations.

\paragraph{Viscosity conditioning.}
By passing the viscosity coefficient $\nu$ as an explicit input (via FiLM modulation), a single TCAN model handles dynamics ranging from near-inviscid shock formation ($\nu \approx 0.001$) to strongly diffusive smooth solutions ($\nu \approx 0.5$) without any retraining. This conditioning is essential for generalization: without it, the model trained at one regime fails qualitatively at another.

\paragraph{FNO benchmark and evaluation.}
Switching to the official FNO dataset~\cite{li2021fourier}, generated by the same pseudo-spectral solver used in that paper, allowed us to compare TCAN directly against published baselines (FNO, RBM, FCN, LNO, PCA+NN) at four spatial resolutions ($s \in \{85, 141, 211, 421\}$). At the coarsest resolution ($s=85$), TCAN achieves a relative $L^2$ error of $0.30$, compared to FNO's $0.011$. The results are consistent across resolutions measured by multiple metrics: MSE, SSIM, PSNR, mass conservation, energy monotonicity, and PDE residual.

\paragraph{2D extension.}
We generalised the TCAN architecture to 2D by replacing all 1D spatial operations with their 2D counterparts, retaining FiLM viscosity conditioning, causal temporal attention, and the bounded residual correction. The model trains stably on the 2D Burgers dataset, but the spatial predictions remain qualitatively poor at this stage, motivating the targeted improvements described below.

%==============================================================================
\section{Limitations}
\label{sec:conclusion_limitations}
%==============================================================================

\paragraph{Autoregressive error accumulation.}
The most significant limitation of TCAN relative to FNO is the autoregressive rollout itself. FNO is a \emph{single-step operator} that maps the initial condition directly to $u(\cdot, T)$ in one forward pass. TCAN, by contrast, calls the controller once per time step over $T-1$ steps, allowing small prediction errors to compound exponentially. At finer grids ($s=421$), where more rollout steps are needed per unit time, the relative $L^2$ error reaches $0.94$. Scheduled training on the model's own prediction errors (as in RNO~\cite{lippe2023pderefiner}) is the most direct path to reducing this gap.

\paragraph{Resolution scaling.}
Error grows monotonically with spatial resolution: the same architecture that achieves $0.30$ at $s=85$ reaches $0.94$ at $s=421$. The finer grid requires more time steps in the rollout, and each extra step adds another compounding round of prediction error. Resolution-adaptive training or spectral regularisation could mitigate this effect.

\paragraph{Boundary artifacts.}
The current implementation uses symmetric (reflect) padding in all convolutional layers. For Burgers with periodic boundary conditions, this is physically incorrect: the reflected "mirror image" introduces spurious correlations at the domain edges that propagate inward during the rollout. Replacing reflect padding with circular padding would eliminate this artefact at a negligible implementation cost.

\paragraph{2D model maturity.}
The 2D TCAN is a prototype: it trains without divergence but does not yet produce accurate spatial reconstructions. The dataset is smaller, the boundary padding issue is more severe in 2D, and the optimisation landscape is harder. These are tractable engineering problems rather than fundamental barriers.

%==============================================================================
\section{Perspectives}
\label{sec:conclusion_perspectives}
%==============================================================================

\paragraph{Short-term fixes (1D and 2D).}
The most impactful near-term improvements are: \emph{(i)} replace reflect padding with circular padding; \emph{(ii)} apply progressive rollout training (expose the model to its own errors from the first epoch); and \emph{(iii)} increase the 2D dataset size to match the 1D FNO setup (1000 train / 200 test trajectories per resolution). Together these three changes are expected to close a substantial portion of the FNO gap and lift the 2D model to quantitative accuracy.

\paragraph{Post-processing and physics-based correction.}
PDE-Refiner~\cite{lippe2023pderefiner} can recover sharpness lost at shocks by iterative denoising refinement, applied as a post-processing step on top of TCAN predictions. PhysicsCorrect~\cite{huang2025physicscorrect} offers a training-free projection that reduces PDE residual errors without modifying the architecture. Both approaches are drop-in enhancements compatible with the current TCAN.

\paragraph{Extension to incompressible Navier-Stokes.}
The long-term scientific goal is to simulate 2D incompressible Navier-Stokes:
\begin{align}
  \frac{\partial \mathbf{u}}{\partial t} + (\mathbf{u} \cdot \nabla)\mathbf{u}
  &= -\nabla p + \nu \nabla^2 \mathbf{u}, \\
  \nabla \cdot \mathbf{u} &= 0.
\end{align}
Enforcing the divergence-free constraint $\nabla \cdot \mathbf{u} = 0$ at each rollout step requires an additional Hodge-projection layer (or a learned pressure-correction step) that is absent from the current Burgers architecture. With this addition, the NFTM framework becomes applicable to the full range of incompressible flows, unlocking realistic engineering applications in aerodynamics and heat transfer.

\paragraph{Data efficiency.}
For Navier-Stokes, generating high-fidelity 2D training trajectories is computationally expensive. Active learning~\cite{musekamp2025activelearning}, which intelligently selects the most informative initial conditions and viscosity values for labelling, is a principled way to cover the dynamical parameter space with fewer than 1000 trajectories, and is a natural next step once the 2D Burgers model is stabilised.

%==============================================================================
\section{Key Takeaway}
\label{sec:conclusion_takeaway}
%==============================================================================

This project demonstrates that the NFTM framework with a TCAN controller successfully learns to simulate 1D Burgers dynamics in a data-driven, physics-aware manner. TCAN outperforms the CNN and Transformer baselines on long-horizon rollouts, and viscosity conditioning via FiLM enables a single model to operate across the full range of flow regimes, from sharp shock formation at $\nu \approx 0.001$ to smooth diffusive behaviour at $\nu \approx 0.5$, without retraining.

The performance gap with respect to FNO ($0.30$ vs.\ $0.011$ relative $L^2$ at $s=85$) is not a fundamental architectural limitation but a consequence of the autoregressive formulation: TCAN accumulates small per-step errors over 256 rollout steps, whereas FNO produces a direct one-shot prediction. The identified fixes (circular boundary padding, progressive rollout training, and a larger 2D dataset) are tractable engineering improvements that are expected to close a substantial portion of this gap.

More broadly, this work establishes a solid foundation for generalizable neural PDE solvers. The NFTM framework, with its modular separation of controller and external memory, is well-positioned to scale to 2D incompressible Navier-Stokes once the 1D Burgers bottlenecks are resolved, a promising direction for data-driven simulation of realistic fluid dynamics.

